\section{Introduction}\label{sec:introduciton}

UI component libraries such as material-ui, ant-design, and element-plus have become foundational infrastructure for modern front-end development~\cite{delgado2016reusing, abdalkareem2017developers, lazuardy2022modern}. They provide reusable interface foundations and are extensively employed in industrial and open-source ecosystems~\cite{decan2019empirical, naik2023awesome, weeraddana2024dependency}. Large open-source communities maintain popular libraries such as material-ui\footnote{https://mui.com/} and ant-design\footnote{https://ant.design/}, which are widely adopted across both industrial and open-source software ecosystems. Consequently, defects in a single component can affect a large number of downstream applications~\cite{venturini2023depended}. Unlike traditional utility libraries, however, UI components are defined not only by functional correctness but also by rich behavioral contracts governing interactions among component properties, user events, accessibility semantics, rendering behavior, and composition contexts~\cite{barr2014oracle,alshayban2020accessibility, zhou2025declarui, bajammal2021semantic}.

These behavioral contracts are only partially reflected by conventional execution-based coverage metrics. Statement and branch coverage indicate whether component code has been executed, but they provide limited insight into whether tests actually validate the behavioral relations that define component correctness~\cite{inozemtseva2014coverage,gopinath2014code,kochhar2017code,andrews2005mutation, schuler2011assessing}. A test may render a component, trigger user interactions, or execute the same control-flow path while failing to verify that disabled components suppress interaction, keyboard navigation preserves accessibility semantics, or responsive layouts maintain expected visual behavior. Consequently, high execution coverage does not necessarily imply high behavioral validation.

Metamorphic relations (MRs) provide a natural representation of such behavioral expectations~\cite{zhou2018metamorphic, zhang2023automated, li2025metamorphic}. Rather than specifying expected outputs for individual executions, an MR describes how observable behavior should change or remain invariant—under systematic transformations of component properties, interactions, state, or execution context~\cite{chen2018metamorphic,cho2025metamorphic,tsigkanos2023large}. MRs have been widely used to alleviate the oracle problem by generating follow-up tests or constructing relational test oracles in domains such as machine learning, scientific computing, and compiler testing~\cite{tian2018deeptest,zhang2018deeproad,chen2016empirical, NolascoMDGGPUAF24}. UI component APIs similarly imply behavioral relations across property changes, accessibility states, rendering contexts, and user interactions, suggesting that inferred MRs can serve not only as guidance for testing but also as an empirical behavioral reference for assessing existing test suites.

Despite these advances, existing work primarily uses MRs to generate tests or construct test oracles. Comparatively little attention has been paid to using inferred MRs as a behavioral specification for assessing the adequacy of existing UI component test suites. Existing execution-based metrics indicate whether component code is exercised, but provide limited evidence of whether tests explicitly validate the behavioral relations implied by component APIs and documentation~\cite{huo2014improving, goldstein2024property}. Consequently, developers still lack a practical way to measure the behavioral validation of UI component test suites beyond execution coverage.

This paper addresses this gap by presenting an MR-based framework for assessing the behavioral validation of UI component test suites. Rather than generating new tests, our framework uses inferred MRs as a behavioral reference space for evaluating existing tests. Given a component's implementation, public API specification, and test suite, the framework (1) infers a component-level MR space using a UI-specific taxonomy, (2) aligns existing tests with the inferred behavioral relations through hybrid deterministic and semantic analysis, and (3) distinguishes between \textit{behavioral reach} (\textit{Touch}) and \textit{behavioral validation} (\textit{Cover}). This relation-level perspective identifies whether behavioral gaps arise from missing test scenarios or from insufficient behavioral validation in tests that already exercise the relevant behavior.

While prior metamorphic testing work often uses MRs for follow-up test generation or oracle construction, we use inferred MRs as a behavioral reference for assessing the adequacy of existing UI component test suites.
We evaluate the framework on 214 components from four widely used UI component libraries. To assess the reliability of the inferred behavioral reference, we manually evaluate both MR inference and test--MR alignment: the inferred MRs achieve 88.6\% usability on a stratified sample, and hybrid alignment reaches an F1 of 96.6\% for \textit{Touch} and 89.7\% for \textit{Cover}. The results show that existing test suites exercise substantially more inferred behavioral relations than they explicitly validate, indicating that many behavioral relations are exercised but lack explicit behavioral validation. Furthermore, MR coverage provides a complementary perspective beyond statement and branch coverage by capturing behavioral validation not reflected by execution-based metrics. Finally, analyses of reported issue descriptions, an oracle-strengthening study, and an MR-relevant injected fault study support the practical relevance of MR coverage for assessing behavioral validation in existing UI component test suites. 


In summary, the main contributions of this paper are:
\begin{itemize}
    \item \textbf{A UI-specific MR taxonomy.}
    We define a six-category taxonomy of component-level metamorphic relations covering input/prop, state/event semantics, interaction/accessibility, visual/layout, composition/context, and data flow.
    \item \textbf{A hybrid LLM-assisted framework for MR-based UI behavioral validation.}
    The framework combines taxonomy-constrained LLM inference with deterministic local-evidence alignment to determine which inferred behavioral relations are exercised (\textit{Touch}) and which are explicitly validated (\textit{Cover}) by existing tests.
    \item \textbf{An empirical study on 214 UI components.} 
    We evaluate the framework on 214 UI components and show that MR coverage complements execution-based coverage by distinguishing behavioral reach from behavioral validation.
    \item \textbf{Evidence of practical relevance.} Through analyses of reported issue descriptions, an oracle-strengthening study, and an MR-relevant injected fault study, we provide evidence that the inferred MR space captures behaviors reflected in real-world issues and that MR coverage offers a practically meaningful perspective for assessing behavioral validation.
\end{itemize}





 
